\subsection{INVARIANT BILINEAR FORMS}

Let g be a Lie algebra over K. The adjoint representation of g on g and the zero representation of g on K define a g-module structure on the K-module

\(N = \mathcal{L}(g, g; K)\) of bilinear forms on g. Briefly we say that a bilinear form \(\beta\) on g is invariant if it is invariant under the representation \(x \mapsto x_{N}\) . By formula (10) the necessary and sufficient condition for this to be so is that:

\[
\beta ([ x, y ], z) = \beta (x, [ y, z ])\tag{13}
\]

for all \(x, y, z\) in \(\mathfrak{g}\) .

Now let \(\mathfrak{d}\) be the Lie algebra of derivations of \(\mathfrak{g}\) . The identity representation of \(\mathfrak{d}\) and the zero representation of \(\mathfrak{d}\) on \(\mathrm{K}\) define a representation \(\mathrm{D} \mapsto \mathrm{D}_{\mathrm{N}}\) of \(\mathfrak{d}\) on \(\mathrm{N}\) . Briefly we say that a bilinear form on \(\mathfrak{g}\) is completely invariant if it is invariant under the representation \(\mathrm{D} \mapsto \mathrm{D}_{\mathrm{N}}\) . A completely invariant bilinear form is invariant. For a bilinear form \(\beta\) on \(\mathfrak{g}\) to be completely invariant, it is necessary and sufficient that:

\[
\beta (\mathrm{D} x, y) + \beta (x, \mathrm{D} y) = 0\tag{14}
\]

for all \(x, y\) in \(\mathfrak{g}\) and \(\mathrm{D} \in \mathfrak{d}\) .

PROPOSITION 7. Let \(\mathfrak{g}\) be a Lie algebra, \(\beta\) an invariant symmetric bilinear form on \(\mathfrak{g}\) and \(\mathfrak{a}\) an ideal of \(\mathfrak{g}\) .

\begin{enumerate}
\def\labelenumi{(\alph{enumi})}
\item
  The orthogonal \(\mathfrak{a}'\) of \(\mathfrak{a}\) with respect to \(\beta\) is an ideal of \(\mathfrak{g}\) .
\item
  If \(\mathfrak{a}\) is characteristic and \(\beta\) is completely invariant, \(\mathfrak{a}'\) is characteristic.
\item
  If \(\beta\) is non-degenerate, \(\mathfrak{a} \cap \mathfrak{a}'\) is commutative.
\end{enumerate}

Let D be a derivation of g. Suppose that a is stable under D and that \(\beta(\mathrm{D}x,y)+\beta(x,\mathrm{D}y)=0\) for x,y in g. Then \(z\in a'\) implies \(Dz\in a'\) , since, for all \(t\in a\) , \(Dt\in a\) and hence \(\beta(\mathrm{D}z,t)=-\beta(z,\mathrm{D}t)=0\) . Thus \(a'\) is stable under D. This establishes (a) and (b).

Now let \(\mathfrak{b}\) be an ideal of \(\mathfrak{g}\) and suppose that the restriction of \(\beta\) to \(\mathfrak{b}\) is zero. For \(x, y\) in \(\mathfrak{b}\) and \(z \in \mathfrak{g}\) , \(\beta ([ x, y ], z) = \beta(x, [y, z]) = 0\) , for \([y, z] \in \mathfrak{b}\) . Thus \([\mathfrak{b}, \mathfrak{b}]\) is orthogonal to \(\mathfrak{g}\) . If \(\beta\) is non-degenerate, \(\mathfrak{b}\) is therefore commutative. This result applied to \(\mathfrak{a} \cap \mathfrak{a}'\) proves (c).

DEFINITION 4. Let g be a Lie K-algebra and M a g-module. Suppose that M, considered as a K-module, admits a finite basis. The bilinear form associated with the g-module M (or with the corresponding representation) is the symmetric bilinear form \((x,y)\mapsto\mathrm{Tr}(x_{\mathrm{M}} y_{\mathrm{M}})\) on g. If the representation in question is the adjoint representation, the associated bilinear form is called the Killing form of g.

PROPOSITION 8. Let \(\mathfrak{g}\) be a Lie algebra and \(\mathbf{M}\) a \(\mathfrak{g}\) -module. Suppose that \(\mathbf{M}\) , considered as a K-module, admits a finite basis. The bilinear form associated with \(\mathbf{M}\) is invariant. For \(x, y, z\) in \(\mathfrak{g}\) , we have:

\[
\begin{array}{r l} & {\mathrm{Tr} ([ x, y ] _ {\mathrm{M}} z _ {\mathrm{M}}) = \mathrm{Tr} (x _ {\mathrm{M}} y _ {\mathrm{M}} z _ {\mathrm{M}}) - \mathrm{Tr} (y _ {\mathrm{M}} x _ {\mathrm{M}} z _ {\mathrm{M}}) = \mathrm{Tr} (x _ {\mathrm{M}} y _ {\mathrm{M}} z _ {\mathrm{M}}) - \mathrm{Tr} (x _ {\mathrm{M}} z _ {\mathrm{M}} y _ {\mathrm{M}})} \\ & {\quad = \mathrm{Tr} (x _ {\mathrm{M}} [ y, z ] _ {\mathrm{M}}).} \end{array}
\]

PROPOSITION 9. Suppose that K is a field and that the Lie algebra g is finite-dimensional over K. Let a be an ideal of g, β the Killing form of g and β' the Killing form of a. Then β' is the restriction of β to a.

Let u be an endomorphism of the vector space g which leaves a stable. Let v be the restriction of u to a and w the endomorphism of the vector space g/a derived from u when passing to the quotient. Then \(Tr u = Tr v + Tr w\) as is seen by taking a basis \((x_{1}, \ldots, x_{n})\) of g of which the first p elements form a basis of a. Then let \(x \in a, y \in a\) and apply the above formula to the case where \(u = (\mathrm{ad}_{\mathrm{g}} x)(\mathrm{ad}_{\mathrm{g}} y)\) . Then \(v = (\mathrm{ad}_{\mathrm{a}} x)(\mathrm{ad}_{\mathrm{a}} y)\) and w = 0. Hence \(\beta(x, y) = \beta'(x, y)\) .

PROPOSITION 10. Suppose that K is a field and that the Lie algebra g is finite-dimensional over K. The Killing form β of g is completely invariant.

Let D be a derivation of g. There exists a Lie algebra \(g'\) containing g as an ideal of codimension 1 and an element \(x_{0}\) of \(g'\) such that \(Dx = [x_{0}, x]\) for all \(x \in g\) (§ 1, no. 8, Example 1). Let \(\beta'\) be the Killing form of \(g'\) . For x, y in g, \(\beta'([x, x_{0}], y) = \beta'(x, [x_{0}, y])\) , that is \(\beta'(Dx, y) + \beta'(x, Dy) = 0\) . Now the restriction of \(\beta'\) to g is \(\beta\) (Proposition 9). Hence the proposition.

